\documentclass[11pt]{article}
\usepackage[a4paper,margin=26mm]{geometry}
\usepackage{amsmath,amssymb,mathtools}
\usepackage{unicode-math}
\setmainfont{Latin Modern Roman}
\setmathfont{Latin Modern Math}
\usepackage{xurl,hyperref}
\hypersetup{hidelinks,pdftitle={UnsolvedMath Problem 10300019}}
\newcommand{\sref}[2]{\hyperlink{source:#1:#2}{[S#2]}}
\newcommand{\cataloguescope}[1]{#1}
\begin{document}
% UnsolvedMath ID 10300019; source catalogue code AMR-102-0019
\setcounter{section}{10300018}
\section{Branched surfaces and triangulations — Question 7.4}\label{10300019}
{\small\sffamily UnsolvedMath ID 10300019 · Topology}
\subsection{Definitions and mathematical statement}

\paragraph{Shared notation.}$\mathbb N=\{1,2,\ldots\}$, and $\mathbb Z,\mathbb Q,\mathbb R,\mathbb C$ denote the integers, rationals, reals and complex numbers. Any other conventions are specified locally.


A surface lamination $\Lambda$ is a closed subset of $M$ locally homeomorphic to a product of a disk with a closed transverse set. It is essential when it has no sphere leaf or torus leaf bounding a solid torus, and its complementary regions are irreducible, have incompressible boundary, and have no compressing monogon. Such a monogon is a compressing disk with one boundary arc on a boundary leaf and the other along a transverse interval fiber. It is genuine if some complementary region is not an interval bundle. A sublamination is a closed union of leaves. Relative to a triangulation, a normal surface or lamination meets each tetrahedron in triangle and quadrilateral disk types; quadrilateral types must be compatible. The Haken sum for compatible normal surfaces is the cut-and-reglue operation along intersections that adds their normal coordinates.
\paragraph{Statement recorded in the supplied source.}
For two laminations in normal form for a fixed triangulation, determine when a well-defined Haken sum exists and how to perform it. The problem is not restricted to transversely measured laminations. The source points to compatibility of quadrilateral types, and more generally simultaneous carrying by a branched surface, as prerequisites to investigate. \sref{10300019}{2}
\subsection{Short English statement}
Determine when the normal-surface cut-and-reglue sum extends to a pair of normal laminations without assuming transverse measures.
\subsection{Sources}
\begin{description}
\item[\hypertarget{source:10300019:1}{S1}] ULAM.AI and the UnsolvedMath Contributors, \emph{UnsolvedMath: A Curated Collection of Open Mathematics Problems}, record 10300019 (AMR-102-0019). CC BY 4.0. \url{https://huggingface.co/datasets/ulamai/UnsolvedMath}.
\item[\hypertarget{source:10300019:2}{S2}] Danny Calegari, \emph{Problems in foliations and laminations of 3-manifolds}, version 0.78 (2002), Question 7.4; Definitions 1.1--1.10 and local definitions and remarks. \url{https://arxiv.org/abs/math/0209081v1}.
\end{description}
\label{end:10300019}
\end{document}
